% ============================================================
\section{Measure Theory}
% ============================================================

\subsection{$\sigma$-Algebras and Measures}

\begin{definition}{$\sigma$-Algebra}{sigma-alg}
    Let $X$ be a non-empty set. A collection $\mathcal{M} \subseteq \mathcal{P}(X)$
    is called a \emph{$\sigma$-algebra} on $X$ if
    \begin{enumerate}
        \item $X \in \mathcal{M}$,
        \item $A \in \mathcal{M} \Rightarrow A^c \in \mathcal{M}$,
        \item $\{A_n\}_{n=1}^{\infty} \subseteq \mathcal{M}
              \Rightarrow \bigcup_{n=1}^{\infty} A_n \in \mathcal{M}$.
    \end{enumerate}
    The pair $(X, \mathcal{M})$ is called a \emph{measurable space}.
\end{definition}

\begin{definition}{Measure}{measure}
    A function $\mu : \mathcal{M} \to [0,+\infty]$ is a \emph{measure} on
    $(X,\mathcal{M})$ if
    \begin{enumerate}
        \item $\mu(\emptyset) = 0$, and
        \item (\emph{countable additivity}) for every sequence of pairwise
              disjoint sets $\{A_n\} \subseteq \mathcal{M}$,
              \[
                  \mu\!\left(\bigsqcup_{n=1}^{\infty} A_n\right)
                  = \sum_{n=1}^{\infty} \mu(A_n).
              \]
    \end{enumerate}
    The triple $(X, \mathcal{M}, \mu)$ is a \emph{measure space}.
\end{definition}

\mn{Borel $\sigma$-algebra: generated by open sets of a topology.}

\begin{theorem}{Continuity of Measure}{cont-meas}
    Let $(X,\mathcal{M},\mu)$ be a measure space.
    \begin{enumerate}
        \item \textup{(Upward continuity)}  If $A_1 \subseteq A_2 \subseteq \cdots$
              in $\mathcal{M}$, then
              $\mu\!\bigl(\bigcup_n A_n\bigr) = \lim_{n\to\infty}\mu(A_n)$.
        \item \textup{(Downward continuity)}  If $A_1 \supseteq A_2 \supseteq \cdots$
              in $\mathcal{M}$ and $\mu(A_1)<\infty$, then
              $\mu\!\bigl(\bigcap_n A_n\bigr) = \lim_{n\to\infty}\mu(A_n)$.
    \end{enumerate}
\end{theorem}

\begin{proof}
    \textbf{(1)} Define $B_1 = A_1$ and $B_n = A_n \setminus A_{n-1}$ for
    $n \geq 2$. The sets $\{B_n\}$ are pairwise disjoint and
    $\bigcup_{k=1}^{n} B_k = A_n$, so by countable additivity
    \[
        \mu\!\left(\bigcup_{n=1}^{\infty} A_n\right)
        = \mu\!\left(\bigsqcup_{n=1}^{\infty} B_n\right)
        = \sum_{n=1}^{\infty} \mu(B_n)
        = \lim_{n\to\infty}\sum_{k=1}^{n}\mu(B_k)
        = \lim_{n\to\infty}\mu(A_n).
    \]

    \textbf{(2)} Apply part (1) to $C_n = A_1 \setminus A_n$, noting
    $\mu(A_1) < \infty$.
\end{proof}

% ============================================================
\section{The Lebesgue Integral}
% ============================================================

\subsection{Construction}

\begin{definition}{Measurable Function}{meas-fn}
    Let $(X,\mathcal{M})$ and $(Y,\mathcal{N})$ be measurable spaces.
    A function $f: X \to Y$ is \emph{measurable} if
    $f^{-1}(B) \in \mathcal{M}$ for every $B \in \mathcal{N}$.
\end{definition}

\begin{lemma}{Monotone Convergence Theorem}{mct}
    Let $\{f_n\}$ be a sequence of non-negative measurable functions with
    $f_n \nearrow f$ pointwise. Then
    \[
        \lim_{n\to\infty} \int f_n \dmu = \int f \dmu.
    \]
\end{lemma}

\begin{theorem}{Dominated Convergence Theorem}{dct}
    Let $\{f_n\}$ be measurable functions with $f_n \to f$ $\mu$-a.e.
    If there exists $g \in L^1(\mu)$ such that $\abs{f_n} \leq g$ $\mu$-a.e.\
    for all $n$, then $f \in L^1(\mu)$ and
    \[
        \lim_{n\to\infty} \int f_n \dmu = \int f \dmu.
    \]
\end{theorem}

\begin{proof}
    Since $g \pm f_n \geq 0$ and $g \pm f_n \to g \pm f$ $\mu$-a.e.,
    Fatou's lemma gives
    \[
        \int (g \pm f)\dmu \leq \liminf_{n\to\infty}\int(g\pm f_n)\dmu.
    \]
    Subtracting $\int g\,\dmu < \infty$ and combining the two inequalities yields
    $\limsup \int f_n\dmu \leq \int f\dmu \leq \liminf \int f_n\dmu$.
\end{proof}

\begin{corollary}{}{dct-cor}
    Under the hypotheses of \cref{thm:dct},
    $\norm{f_n - f}_{L^1} \to 0$ as $n \to \infty$.
\end{corollary}

% ============================================================
\section{\texorpdfstring{$L^p$}{Lp} Spaces}
% ============================================================

\begin{definition}{$L^p$ Space}{lp-space}
    For $1 \leq p < \infty$ define
    \[
        L^p(\mu) \;=\; \Bigl\{\, f \text{ measurable} :
            \norm{f}_{p} \coloneqq \Bigl(\int \abs{f}^p\dmu\Bigr)^{1/p} < \infty \,\Bigr\}.
    \]
    For $p=\infty$,\; $\norm{f}_{\infty} \coloneqq \esssup\abs{f}$.
\end{definition}

\begin{theorem}{Hölder's Inequality}{holder}
    Let $p,q\in(1,\infty)$ be conjugate exponents, $\tfrac{1}{p}+\tfrac{1}{q}=1$.
    For $f\in L^p$, $g\in L^q$,
    \[
        \int \abs{fg}\dmu \;\leq\; \norm{f}_{p}\,\norm{g}_{q}.
    \]
\end{theorem}

\begin{example}{Convolution on $\R^n$}{conv}
    For $f\in L^1(\R^n)$ and $g\in L^p(\R^n)$ the convolution
    \[
        (f*g)(x) = \int_{\R^n} f(x-y)\,g(y)\dx[y]
    \]
    satisfies $\norm{f*g}_{p} \leq \norm{f}_{1}\norm{g}_{p}$ (Young's inequality).
\end{example}

\begin{remark}{On completeness}{completeness}
    Fischer--Riesz: $L^p(\mu)$ is a Banach space for every $1\leq p\leq\infty$.
    For $p=2$ it is a Hilbert space with inner product
    $\inner{f}{g} = \int f\,\ol{g}\dmu$.
\end{remark}

% ============================================================
%  Quick-reference table
% ============================================================
\section*{Quick-Reference: Key Inequalities}
\addcontentsline{toc}{section}{Quick-Reference: Key Inequalities}

\begin{center}
\renewcommand{\arraystretch}{1.4}
\begin{tabular}{@{} l l l @{}}
    \toprule
    \textsf{Name} & \textsf{Statement} & \textsf{Conditions} \\
    \midrule
    Markov     & $\mu(\abs{f}\geq\lambda)\leq \lambda^{-1}\norm{f}_1$
               & $f\in L^1,\;\lambda>0$ \\
    Chebyshev  & $\mu(\abs{f}\geq\lambda)\leq \lambda^{-2}\norm{f}_2^2$
               & $f\in L^2,\;\lambda>0$ \\
    Hölder     & $\norm{fg}_1 \leq \norm{f}_p\norm{g}_q$
               & $\tfrac1p+\tfrac1q=1$ \\
    Minkowski  & $\norm{f+g}_p \leq \norm{f}_p+\norm{g}_p$
               & $1\leq p\leq\infty$ \\
    Jensen     & $\phi\!\bigl(\int f\dmu\bigr)\leq\int\phi\circ f\dmu$
               & $\phi$ convex, $\mu$ a prob.\ measure \\
    \bottomrule
\end{tabular}
\end{center}
